\section{Gobernanza y FSM Determinista: Estados y Reglas de Transición}
\label{sec:fsm}

\subsection{Definición Formal de la FSM}
La máquina de estados finitos que gobierna la negociación se define formalmente como una tupla:
\begin{equation}
    \mathcal{M} = (S, \Sigma, \delta, s_0, F)
\end{equation}
donde:
\begin{itemize}
    \item $S = \{\texttt{on\_board}, \texttt{board\_format\_err}, \texttt{validation}, \texttt{val\_format\_err}, \texttt{ended}, \texttt{frozen}\}$ representa el conjunto de estados internos admisibles.
    \item $\Sigma$ representa el alfabeto de eventos y acciones entrantes (\texttt{propose}, \texttt{counter\_propose}, \texttt{validate}, \texttt{confirm}, \texttt{validation\_rejected}, \texttt{report\_format\_error}, \texttt{fix\_format}, \texttt{freeze}).
    \item $\delta: S \times \Sigma \rightarrow S$ es la función de transición determinista ejecutada proceduralmente en C\#.
    \item $s_0 = \texttt{on\_board}$ es el estado inicial tras completarse con éxito el handshake.
    \item $F = \{\texttt{ended}, \texttt{frozen}\}$ es el conjunto de estados terminales inmutables.
\end{itemize}

\subsection{Mapeo Bidireccional de Estados}
El estándar público A2A define estados genéricos para las tareas de agentes. El orquestador mantiene un mapeo bidireccional estricto entre estos estados y los estados del motor interno de C\#, tal como se sintetiza en el Cuadro~\ref{tab:fsm_mapping}.

\begin{table}[htbp]
\centering
\small
\begin{tabular}{lllp{6.5cm}}
\toprule
\textbf{Estado A2A (Público)} & \textbf{Estado FSM (C\#)} & \textbf{Naturaleza} & \textbf{Significado y Comportamiento} \\
\midrule
\texttt{working} & \texttt{on\_board} & Transitorio & Fase activa de negociación, diálogo persuasivo e intercambio de contraofertas. \\
\texttt{input\_required} & \texttt{board\_format\_err} & Recuperable & Payload inválido sintácticamente; activa el bucle de autocorrección cognitiva. \\
\texttt{evaluating} & \texttt{validation} & Auditoría & Evaluación de la propuesta consolidada frente al Deterministic Guard. \\
\texttt{input\_required} & \texttt{val\_format\_err} & Recuperable & Discrepancia de esquema de datos entre el Fact Sheet y el backend de negocio. \\
\texttt{completed} & \texttt{ended} & Terminal & Acuerdo formalmente alcanzado, validado por ambas partes y firmado. \\
\texttt{failed} & \texttt{frozen} & Terminal / Pausa & Negociación congelada tras superar el límite de reintentos ($N > 3$) o rechazo definitivo. \\
\bottomrule
\end{tabular}
\caption{Mapeo Bidireccional entre el estándar A2A y la FSM interna del Orquestador.}
\label{tab:fsm_mapping}
\end{table}

\subsection{Matriz de Transiciones Legales de la FSM}
El motor procedimental en C\# (\texttt{NegotiationFsmEngine}) aplica de forma estricta la función de transición $\delta$, rechazando cualquier acción fuera de matriz mediante \texttt{InvalidTransitionException}, tal como se detalla en el Cuadro~\ref{tab:fsm_transition_matrix}.

\begin{table}[htbp]
\centering
\small
\begin{tabular}{lp{4.2cm}lp{4.8cm}}
\toprule
\textbf{Estado Origen ($s_t$)} & \textbf{Acción / Evento ($a$)} & \textbf{Estado Destino ($s_{t+1}$)} & \textbf{Efecto en Turno y Persistencia} \\
\midrule
\texttt{on\_board} & \texttt{propose}, \texttt{counter\_propose} & \texttt{on\_board} & Alterna turno al contraparte; persiste facts. \\
\texttt{on\_board} & \texttt{validate} & \texttt{validation} & Transfiere a auditoría de reglas duras. \\
\texttt{on\_board} & \texttt{report\_format\_error} & \texttt{board\_format\_err} & Mantiene turno en emisor para autocorrección. \\
\texttt{on\_board} & \texttt{accept}, \texttt{complete} & \texttt{ended} & Cierra sesión; libera turno (\texttt{activeTurn} nulo). \\
\texttt{on\_board} & \texttt{freeze}, \texttt{reject}, \texttt{cancel} & \texttt{frozen} & Estado terminal/pausa; bloquea turnos futuros. \\
\midrule
\texttt{board\_format\_err} & \texttt{fix\_format}, \texttt{retry} & \texttt{on\_board} & Autocorrección exitosa; reanuda diálogo. \\
\texttt{board\_format\_err} & \texttt{freeze}, \texttt{abort} & \texttt{frozen} & Superación de reintentos o fallo irrecuperable. \\
\midrule
\texttt{validation} & \texttt{confirm}, \texttt{validation\_passed} & \texttt{ended} & Acuerdo auditado y firmado con éxito. \\
\texttt{validation} & \texttt{validation\_rejected}, \texttt{renegotiate} & \texttt{on\_board} & Rechazo por guard; retorna a diálogo con feedback. \\
\texttt{validation} & \texttt{report\_format\_error} & \texttt{val\_format\_err} & Discrepancia de esquema de backend en validación. \\
\texttt{validation} & \texttt{freeze}, \texttt{reject} & \texttt{frozen} & Cancelación durante fase de validación. \\
\midrule
\texttt{val\_format\_err} & \texttt{fix\_format}, \texttt{validate} & \texttt{validation} & Subsanación de esquema; reintenta validación. \\
\texttt{val\_format\_err} & \texttt{freeze}, \texttt{abort} & \texttt{frozen} & Abandono tras error de esquema. \\
\midrule
\texttt{ended} / \texttt{frozen} & \textit{Cualquier acción} & \textit{Rechazado} & Lanza \texttt{NegotiationTerminalException}. \\
\bottomrule
\end{tabular}
\caption{Matriz de Transiciones Legales ejecutada por \texttt{NegotiationFsmEngine}.}
\label{tab:fsm_transition_matrix}
\end{table}

\subsection{Ciclo de Negociación y Dinámica de Transiciones}
La interacción dinámica entre los agentes negociadores se rige por una secuencia estricta de fases operativas y salvaguardas procedurales, tal como se ilustra en la Figura~\ref{fig:negotiation_cycle}. El ciclo se estructura en tres bloques fundamentales:

\begin{enumerate}
    \item \textbf{Bloque 1: Negociación Activa (\texttt{on\_board}, \texttt{board\_format\_err}, \texttt{frozen}):} Constituye el diálogo comercial recurrente. El agente proponente transmite una oferta inicial en estado \texttt{on\_board}. La contraparte puede responder aceptando o emitiendo una contraoferta en el mismo estado. Si se detecta un error de sintaxis o de esquema en el payload, el orquestador retorna \texttt{board\_format\_err}, obligando al agente a subsanar el payload e intentar de nuevo. Si concurren condiciones de espera externa o bloqueo transitorio, la negociación se suspende temporalmente en \texttt{frozen} hasta su reanudación controlada.
    \item \textbf{Bloque 2: Validación de Negocio (\texttt{validation}, \texttt{val\_format\_err}):} Al converger hacia un acuerdo preliminar satisfactorio, se emite una solicitud en estado \texttt{validation} para someter el \textit{Fact Sheet} al \textit{Deterministic Guard}. Si la estructura del mensaje de validación es errónea, se activa el estado \texttt{val\_format\_err} para corrección. Si las reglas duras de negocio (márgenes comerciales, disponibilidad de inventario o límites de crédito) son infringidas, el orquestador rechaza el cierre y retrotrae la negociación al estado \texttt{on\_board} para su reformulación.
    \item \textbf{Bloque 3: Finalización Contractual (\texttt{ended}):} Una vez que ambas partes han validado satisfactoriamente los términos cuantitativos y legales frente a sus respectivos backends, la FSM transiciona de forma irreversible al estado terminal \texttt{ended}, procediendo a la firma del acuerdo y al archivado de la bitácora en la persistencia \textit{Dual-Track}.
\end{enumerate}

\begin{figure}[htbp]
    \centering
    \includegraphics[width=0.75\textwidth,height=0.78\textheight,keepaspectratio]{negotiation_cycle.png}
    \caption{Ciclo de negociación gobernado por la FSM determinista: alternancia de turnos y diálogo en \texttt{on\_board}, bucle de autocorrección sintáctica (\texttt{board\_format\_err}), suspensión temporal (\texttt{frozen}), compuerta de validación contra reglas de negocio (\texttt{validation}, \texttt{val\_format\_err}) y formalización final (\texttt{ended}).}
    \label{fig:negotiation_cycle}
\end{figure}

\subsection{Persistencia Dual-Track: Abstracciones y Puertos de Dominio}
Para preservar la naturaleza 100\% \textit{stateless} del Orquestador sin comprometer la latencia ni la trazabilidad legal, el diseño desacopla el almacenamiento en dos puertos independientes bajo el principio de Inversión de Dependencias (DIP):

\begin{enumerate}
    \item \textbf{Puerto de Estado Vivo (\texttt{IHotStateStore} / Track 2):} Almacén de lectura/escritura síncrono de latencia sub-milisegundo ($<2$~ms). Rehidrata en cada turno el snapshot \texttt{HotStateSession} (\texttt{NegotiationId}, \texttt{CurrentState}, \texttt{ActiveTurn}, \texttt{LastSenderId}, \texttt{TurnIndex}, \texttt{ConsolidatedFacts}, \texttt{UpdatedAt}) y garantiza mutaciones atómicas.
    \item \textbf{Puerto de Bitácora Forense (\texttt{ITranscriptStore} / Track 1):} Registro inmutable y append-only para auditoría contractual y análisis legal posterior. Su ejecución es desacoplada (out-of-band) y la FSM jamás lee de este almacén para tomar decisiones de transición.
\end{enumerate}

\subsection{Modelo de Excepciones de Dominio y Control de Turnos}
El motor procedural de C\# intercepta y rechaza cualquier comportamiento anómalo o probabilístico mediante excepciones de dominio fuertemente tipadas:
\begin{itemize}
    \item \texttt{TurnConflictException}: Lanzada cuando un agente intenta emitir un mensaje fuera de su turno asignado ($s_{\text{emisor}} \neq \texttt{ActiveTurn}$).
    \item \texttt{TurnOutOfOrderException}: Lanzada cuando el índice de turno recibido no es estrictamente creciente respecto a la secuencia monótona de la sesión ($t_{\text{recibido}} \le t_{\text{actual}}$).
    \item \texttt{InvalidTransitionException}: Lanzada ante acciones ilegales según la matriz $\delta(s_t, a)$.
    \item \texttt{NegotiationTerminalException}: Lanzada categóricamente ante cualquier intento de interactuar sobre una sesión que ya ha alcanzado un estado terminal (\texttt{ended} o \texttt{frozen}).
\end{itemize}

\subsection{Deterministic Guard y Reglas de Negocio}
El componente \textit{Deterministic Guard} actúa como una compuerta infranqueable ante cualquier acuerdo propuesto por los agentes. Su evaluación es binaria y determinista, evaluando las siguientes restricciones operativas:
\begin{enumerate}
    \item \textbf{Margen Mínimo Comercial:} Comprueba que el precio ofertado $P_{\text{acuerdo}}$ respete el umbral mínimo de rentabilidad según el costo base $C_{\text{base}}$:
    \begin{equation}
        \frac{P_{\text{acuerdo}} - C_{\text{base}}}{P_{\text{acuerdo}}} \ge \mu_{\min}
    \end{equation}
    \item \textbf{Disponibilidad de Inventario:} Verifica contra el sistema de gestión (ERP) que la cantidad demandada $Q$ no exceda el stock físico comprometible $S_{\text{disp}}$.
    \item \textbf{Límite de Riesgo de Crédito:} Evalúa la exposición crediticia del tenant solicitante frente a su póliza vigente.
    \item \textbf{Factibilidad de Plazo:} Comprueba que la fecha límite de entrega $T_{\text{entrega}}$ sea compatible con el tiempo de transporte mínimo $T_{\text{lead}}$.
\end{enumerate}

Si cualquiera de las condiciones anteriores resulta insatisfecha, la transición hacia \texttt{ended} es abortada y la FSM retorna a \texttt{on\_board} adjuntando un diagnóstico cuantitativo para renegociación.

\subsection{Bucle de Autocorrección Cognitiva (Self-Correction Loop)}
Cuando un mensaje recibido adolece de errores de sintaxis, no respeta el esquema JSON esperado o intenta un salto de estado ilegal según la matriz $\delta$, el Gateway responde con un código de error HTTP 422 (\textit{Unprocessable Entity}).

El agente cognoscitivo receptor debe implementar el bucle formal de autocorrección:
\begin{enumerate}
    \item Captura del objeto de diagnóstico retornado por el orquestador.
    \item Inyección del mensaje de error en el contexto conversacional del LLM.
    \item Regeneración estricta de la propuesta mediante Structured Outputs.
    \item Limitación imperativa: el contador de reintentos acumulados $n$ está acotado por $n \le N_{\max}$ (con $N_{\max} = 3$). En caso de alcanzarse $n > N_{\max}$, el orquestador transiciona forzosamente la sesión al estado \texttt{frozen} (\texttt{failed}) y deriva la incidencia para mediación humana.
\end{enumerate}

\subsection{Evaluador Multicriterio MCDA para Negociaciones Concurrentes}
En escenarios multilaterales (subastas con $N$ proveedores concurrentes), el Meta-Orquestador resuelve la selección del proveedor adjudicatario empleando un modelo MCDA (\textit{Multi-Criteria Decision Analysis}) basado en una función de utilidad aditiva ponderada:
\begin{equation}
    \text{Score}_k = w_p \cdot \widetilde{P}_k + w_t \cdot \widetilde{T}_k + w_d \cdot \widetilde{D}_k + w_r \cdot \widetilde{R}_k
\end{equation}
donde $\sum w_i = 1$, y los términos representan valores normalizados en el intervalo $[0, 1]$ correspondientes a Precio ($P$), Tiempo de entrega ($T$), Descuentos por volumen ($D$) y Reputación histórica del proveedor ($R$).
El ganador recibe la notificación \texttt{complete\_task}, mientras que los proveedores no seleccionados reciben \texttt{cancel\_task}.
